\documentclass[11pt]{article}
\usepackage[T1]{fontenc}
\usepackage[utf8]{inputenc}
\usepackage{lmodern}
\usepackage{geometry}
\geometry{margin=1in}
\usepackage{hyperref}
\usepackage{enumitem}
\usepackage{xcolor}
\usepackage{listings}
\lstset{basicstyle=\ttfamily\small,breaklines=true,frame=single,columns=fullflexible}
\setlength{\emergencystretch}{3em}

\newcommand{\code}[1]{\texttt{#1}}
\newcommand{\CeTTa}{\textup{CeTTa}}
\newcommand{\HE}{\textup{HE}}
\newcommand{\PeTTa}{\textup{PeTTa}}
\newcommand{\PathMap}{\textup{PathMap}}
\newcommand{\MORK}{\textup{MORK}}

\title{\CeTTa{} Implementation Guide\\
  \large What to Do Next, and Why That Order}
\author{Zar Goertzel \and Oru\v{z}i}
\date{2026-04-11}

\begin{document}
\maketitle

\begin{abstract}
The roadmap describes the territory.  This guide describes the natural
order for walking through it.  Universal term sharing comes first because
it changes what identity means.  Tabling comes second because it builds
on canonical identity.  Everything else follows from those foundations.
This is guidance, not a contract.
\end{abstract}

\tableofcontents

\section{Where We Are}

\CeTTa{} \code{main} already has:

\begin{itemize}
\item \code{TermUniverse} and persistent hash-consing
\item \code{SearchContext} and \code{ChoicePoint} for nondeterminism
\item \code{TableStore} for revision-aware memoization
\item \code{TAIL\_REENTER} for stack-safe deterministic recursion
\item \code{SpaceMatchBackend} for backend abstraction
\end{itemize}

These are real seams, not aspirations.  The gap is that \code{Space}
still owns \code{Atom **atoms} directly.  Spaces are not yet clean
membership layers over a shared universe.

\section{The Natural Order}

\subsection{First: Universal Term Sharing}

Why first?  Because it changes the meaning of identity, equality,
deduplication, and index keys.  Everything downstream---tabling,
backend parity, hot-path work---depends on terms having honest
canonical identity.

The move is simple in concept:
\begin{enumerate}
\item All persistent terms go through canonical storage.
\item Spaces store membership (which canonical terms they contain),
  not ownership (private copies of terms).
\item Indices key on canonical identity, not ad-hoc pointers.
\end{enumerate}

The algebra: a Space is a \emph{predicate} over the term universe---it
answers ``does this term belong?'' rather than ``here is my private copy.''
Membership and identity are cleanly separated.  The specific representation
(arrays, sets, bitmaps, tries) is implementation detail; the invariant is
that no Space owns term storage.

This follows from the HE spec's treatment of spaces as logical containers,
not physical heaps.  The transition means:
\begin{itemize}
\item Terms have one canonical identity regardless of which spaces contain them.
\item Adding a term to multiple spaces does not duplicate storage.
\item Indices can key on stable identities that survive space operations.
\end{itemize}

\paragraph{Current status (April 2026).}
Phase 1a--c is structurally complete: \code{Space} uses \code{AtomId},
\code{TermUniverse} owns persistent terms.  However, capacity witnesses show
100--800$\times$ regression vs.\ pre-TU baselines.  The architecture is correct;
the implementation pays $O(n)$ per store where SOTA pays $O(\mathrm{arity})$.
Phase 1d (byte-packing) recovers $\sim$2$\times$; Phase 1f (AtomId-native
construction) is required to recover the remaining 50--400$\times$.

\paragraph{Council note.}
\emph{Schulz}: ``Perfect term sharing is the converged design.  But sharing
with $O(n)$ per store is sharing done wrong.  Fix the construction path.''
\emph{Conchon--Filli\^{a}tre}: ``Smart constructors enforce sharing at
construction boundaries.  You never hash a term whose children aren't
already interned.''
\emph{Ritchie}: ``Big rewrites go through worse-before-better.  Be honest
about which phase you're in.  You're not in `better' yet.''

\subsection{Phase 1d: Canonical Storage ABI}

After Space membership is authoritative via \code{AtomId} (Phase~1a--c), freeze the
physical canonical store format before building more layers.

The \code{TermUniverse} persistent store should become byte-packed:

\begin{lstlisting}[language=C]
typedef struct {
    uint8_t  tag;          // symbol/var/expr/ground
    uint8_t  subtag;       // grounded subtype or var-mode
    uint16_t arity_or_len; // expr arity or short length
    uint32_t sym_or_head;  // SymbolId / cached head / spelling
    uint32_t aux32;        // payload words
    uint32_t hash32;       // structural hash
} CettaTermHdr;
\end{lstlisting}

Payloads:
\begin{itemize}
\item \code{SYMBOL}: no payload; \code{sym\_or\_head = SymbolId}
\item \code{VAR}: payload = \code{uint64\_t var\_id}
\item \code{EXPR}: payload = \code{AtomId child\_ids[arity]}
\item Stable groundeds: inline or short payload
\item Pointer-bearing groundeds: not in portable canonical ABI yet
\end{itemize}

Keep \code{term\_universe\_get\_atom()} working via lazy decode into
\code{decoded\_cache} field.  Add store-level accessors:
\code{tu\_kind()}, \code{tu\_arity()}, \code{tu\_head\_sym()}, \code{tu\_child()}.

\subsection{Phase 1e: Backend Path ABI}

Add separate structural path encoding for \PathMap{}/\MORK{}/interop:

\begin{lstlisting}[language=C]
bool term_encode_pathkey(Arena *dst,
                         const TermUniverse *u,
                         AtomId id,
                         uint8_t **out_bytes,
                         uint32_t *out_len);
\end{lstlisting}

This format is self-contained, versioned, preorder (outer-first for prefix
sharing).  Do \textbf{not} use universe-local child \code{AtomId} bytes as
the bridge format---use recursive structural encoding.

\paragraph{Import discipline.}
Bridge materialization from \PathMap{}/\MORK{} into \CeTTa{} must use
bottom-up \code{backend\_node $\rightarrow$ AtomId} conversion, memoized by
backend node identity (stable handle, canonical byte span, or persistent
node pointer).  Do \textbf{not} materialize via
\code{dump-text $\rightarrow$ parse-text $\rightarrow$ space\_add\_atom\_id}:
that path pays UTF-8 parse cost, rebuilds transient \code{Atom*} trees, and
re-interns structure the backend already has.  Import cost should be
proportional to the number of unique backend nodes, not to repeated
traversals across rows.

\subsection{Phase 1f: Born-Canonical Persistent Construction}

\textbf{Why this is critical:} Phases 1a--e establish the \emph{structural}
foundation for term sharing, but they introduced a 100--800$\times$ performance
regression. The current architecture computes \code{atom\_hash()} and
\code{term\_universe\_atom\_is\_stable()} as $O(n)$ tree walks on every
store---SOTA systems do $O(\mathrm{arity})$ because children are already interned.

The fix is \textbf{bottom-up construction}: the evaluator works with
\code{AtomId} natively, not \code{Atom*}. When constructing a parent term,
children are already \code{AtomId}s with cached hashes, so:

\begin{lstlisting}[language=C]
// Current: O(n) per store
atom_hash(parent);  // recursively hashes all children

// After 1f: O(arity) per store
parent_hash = combine(child1.hash32, child2.hash32, ...);
\end{lstlisting}

This follows the Conchon--Filli\^{a}tre ``smart constructor'' pattern: sharing
is enforced at construction boundaries, and you never hash a term whose
children aren't already interned.

\paragraph{Implementation path.}
\begin{enumerate}
\item \textbf{Constructor substrate}: Add \code{tu\_intern\_*} constructors
  (\code{tu\_intern\_symbol}, \code{tu\_intern\_var}, \code{tu\_intern\_int},
  \code{tu\_intern\_float}, \code{tu\_intern\_bool}, \code{tu\_intern\_string},
  \code{tu\_expr\_from\_ids}).  These return \code{AtomId} directly with
  $O(\mathrm{arity})$ hashing---children must already be interned.
\item \textbf{Persistent ingestion lanes}: Migrate parser and batch-add paths
  to build \code{AtomId}s bottom-up.  The \code{term\_universe\_store\_atom\_id()}
  path should delegate to constructors internally.
\item \textbf{Counters and capacity witnesses}: Add store-level counters
  (\code{decode\_count}, \code{hash\_recompute\_count}) and assert they remain
  zero for stable-atom-only workloads.  Verify on benchmark suite.
\item \textbf{Runtime reconstruction} (if needed): For runtime paths that
  currently rebuild \code{Atom*} trees before re-interning, add ID-native
  reconstruction.  Only migrate paths where profiling shows hot decode traffic.
\end{enumerate}

\paragraph{Explicit non-goals.}
\begin{itemize}
\item \textbf{No whole-evaluator rewrite}: The evaluator continues to use
  \code{Atom*} views for hot-path access.  This phase improves the store
  boundary, not the eval core.
\item \textbf{No Layer 3 migration}: Converting the evaluator to work natively
  with \code{AtomId} handles is a separate, larger effort (Phase 2+).
\item \textbf{No bridge ABI collapse}: The four-ABI split (logical, store,
  eval, backend) remains.  Do not unify store and backend formats.
\end{itemize}

\paragraph{Compatibility theorem.}
For any \code{Atom*} $a$, if $\code{tu\_intern}(a)$ returns \code{AtomId} $i$,
then $\code{tu\_decode}(i)$ is structurally equal to $a$.  This is the round-trip
invariant that protects existing code during migration.

\paragraph{Exit criteria for Phase 1 (term sharing complete).}
\begin{itemize}
\item \code{tu\_tail\_special\_forms} $\leq 2\times$ fc3afa7 baseline (0.65s)
\item \code{tu\_metamath\_stream\_basic} $\leq 2\times$ fc3afa7 baseline (0.20s)
\item \code{tu\_tilepuzzle} passes within 300s envelope
\item Hash computation is $O(\mathrm{arity})$, not $O(n)$
\end{itemize}

Until these criteria are met, Phase 1 is \emph{structurally complete but
performance-regressed}---commit-ready, but not done.

\subsection{Phase 2: ASM Search Seam (Read/Search v1)}

Once Phase 1 exit criteria are met, introduce a small explicit episode
seam between the language surface and the backend ops.

Positive example: a \code{SearchEpisode} wrapping the existing
\code{SearchContext}, consuming a normalized \code{SearchRequest} and a
\code{SearchLangContract}, exposing \code{visit\_query},
\code{visit\_conjunction}, \code{exact\_contains} through a
\code{SearchMachineOps} dispatch table.

Negative example: a first tranche that already carries MM2, ATP
saturation, solver-oracle lanes, or a scheduler tower.  That is a
later epoch.

\paragraph{Files.}
\code{src/search\_episode.h}, \code{src/search\_episode.c},
\code{src/search\_episode\_default.c}.  The existing
\code{src/search\_machine.h} (branch-local \code{SearchContext} and
\code{ChoicePoint}) is left in place and becomes what the default
episode drives.

\paragraph{Axis factorization (data, not behavior).}
\begin{itemize}
\item \emph{lane}: semantic proof/execution family
  (\code{recursive-dependent-proof} in v1; others deferred)
\item \emph{backend}: storage/indexing engine (native / \PathMap{} /
  \MORK{}), chosen by the actual \code{Space}, not by a global preference
\item \emph{fairness}: DFS / BFS / IDDFS / native (semantic only when
  bounded)
\item \emph{table mode}: \code{none} in v1; variant/subsumptive/incremental
  are explicit future modes, not universal law
\item \emph{order}: presentation contract (native / reverse / lex /
  unspecified)
\item \emph{demand}: unbounded / once / limit-k / streaming fold
\end{itemize}

\paragraph{First tranche (behavior-preserving).}
\begin{enumerate}
\item Add the episode / request / visitor types; no algorithms move.
\item Implement \code{search\_episode\_default} by delegating to
  current \code{SearchContext} + fallback conjunction engine + backend
  direct ops when advertised via a small \code{SearchCaps} struct.
\item Pull \code{search-policy} and \code{pragma!}\ table-mode
  interpretation out of \code{eval.c} into a
  \code{search\_request\_from\_surface(...)} builder.
\item Route the streaming surfaces (\code{select}, \code{once},
  \code{collect}, \code{fold}, \code{fold-by-key}) through the visitor
  seam.  The live glue currently sits around
  \code{eval.c:2797}--\code{eval.c:3375}.
\end{enumerate}

\paragraph{Explicit non-goals for Phase 2.}
\begin{itemize}
\item \textbf{No \code{EvalLaneOps} tier} in v1.  Lane families are a
  later epoch; the v1 seam is read/search only.
\item \textbf{No backend details in the seam}: hash/trie/pathmap
  choice, candidate narrowing, join order, row/provenance bookkeeping,
  exact-memo encoding, and GC remain backend-local.
\item \textbf{No language observables below the seam}: per-\code{--lang}
  promises (\HE{}, \PeTTa{}, future langs) normalize to a
  \code{SearchRequest + SearchLangContract} above the seam.
\item \textbf{No new table modes in v1}: \code{table\_mode = none} is
  the only value honored; variant mode arrives via the plugin wrap in
  the next subsection.
\end{itemize}

\subsection{Phase 3: Tabling as First ASM Plugin}

With canonical term identity (Phase 1) and the episode seam (Phase 2)
in place, tabling lands as the first pluggable table mode under the
seam.

Positive example: \code{TableStore} (the existing revision-aware
store) wrapped as a \code{SearchTableOps} implementation, with
\emph{variant} as the sole honored \code{table\_mode} value in v1.

Negative example: today's variant-table behavior silently becoming
the universal law of all future backends and all future lanes.

\code{TableStore} already exists.  The work is:
\begin{enumerate}
\item Variant keys use canonical term identity (enabled by Phase 1).
\item Wrap \code{TableStore} behind a \code{SearchTableOps} vtable;
  register it as a plugin the episode may consult when
  \code{req.table\_mode = variant}.
\item Decide whether revision-aware invalidation is enough, or
  whether recursive workloads need suspension/resumption (Level 3+ on
  the roadmap's tabling ladder, see \code{cetta\_roadmap.tex}
  \S{}``Tabling: The Ladder of Capabilities'').
\end{enumerate}

Let the workloads tell you.  Recursive backward chaining (where
\PeTTa{} currently completes and \CeTTa{} OOMs) will show whether
simple caching suffices or whether full SLG is required.  Additional
modes (subsumptive, incremental, answer-subsuming, monotonic) stay
explicit---never a hidden default.  The goal is performant backward
chaining that matches \PeTTa{}'s time/memory while preserving \HE{}
semantics.

\subsection{Phase 4: Evaluator Control Convergence}

Why fourth?  Because control-path surgery should respect the
representation and tabling contracts already established.

The key constraint: \code{TAIL\_REENTER} is constitutional for deterministic
singleton forms.  Don't replace it without proving the replacement is
both correct and faster.

Frame-based suspension goes where it is actually needed: multi-result
streaming, typed application, bounded consumers---i.e., exactly the
surfaces routed through the Phase~2 visitor seam.  Positive example:
\code{select} / \code{fold} / \code{fold-by-key} calling
\code{episode\_ops->visit\_query} with a \code{SearchVisitCtl} that
may request \code{STOP} once demand is satisfied.  Negative example:
ad-hoc tail-position tricks re-invented inside every streaming
consumer.

\subsection{Phase 5: Backend Parity as Contract Equivalence}

Why fifth?  Because backend work risks baking in premature
assumptions if the semantic center of gravity is still moving.

Once term identity, the episode seam, and tabling are stable, verify
that native, \PathMap{}, \MORK{}, and DAS backends produce answer
streams observationally equivalent under the same
\code{SearchRequest + SearchLangContract}.  Internal realizations
may diverge freely: native may grow WAM-ish frame tuning and custom
exact-membership structures; \PathMap{} may grow counted answer
stores, prefix scans, and algebraic shortcuts; \MORK{}/MM2 may grow
bridge-specific packet contracts.  That is healthy divergence, not
architectural failure.

The parity claim is: for any two backend realizations $B_1, B_2$
satisfying the same \code{SearchLangContract} on shared logical
content, a conforming visitor observes equivalent answer streams
modulo the contract's declared observability (order, multiplicity,
freshening).  Storage differs.  Internal search and table
realizations may differ.  The contract does not.

\subsection{Phase 6: Hot Paths and Specialized Lanes}

Why last?  Because their honest shape depends on what the earlier work
made cheap or unnecessary.

After universal sharing, tabling, control, and backend parity---then
look at what's still slow.  That's the real hot path, not the one you
imagined before the foundations were in place.

PLN, ATP, factor-graph reasoning: these are specialized lanes over the
same semantic contract, not modifications to the core evaluator.

\section{Dependencies}

The order above reflects dependencies:

\begin{center}
\begin{tabular}{lcl}
Universal term sharing (1a--c) & $\rightarrow$ & Canonical store ABI (1d) \\
Canonical store ABI (1d) & $\rightarrow$ & Backend path ABI (1e) \\
Backend path ABI (1e) & $\rightarrow$ & Tabling keys canonical \\
Canonical tabling & $\rightarrow$ & Control can rely on table contracts \\
Stable control & $\rightarrow$ & Backend parity testable \\
All foundations & $\rightarrow$ & Hot paths are real, not imagined
\end{tabular}
\end{center}

You can do preparatory work in parallel, but the integration order
matters.

\section{What Tests Already Exist}

Use these as sanity checks:
\begin{itemize}
\item \code{tests/test\_canon\_roundtrip\_contracts.metta} --- canonicalization
\item \code{tests/test\_alpha\_equiv.metta} --- alpha equivalence
\item \code{tests/test\_variant\_table\_contracts.metta} --- tabling
\item Backward chaining benchmarks (metamath, PLN) --- performance witnesses
  where \PeTTa{} sets the target
\item \code{make test} --- full regression
\end{itemize}

If these break, something's wrong.  If they pass but behavior is
different, the tests need expansion.

\section{Anytime Work}

Some tasks can proceed in parallel with the ordered work above.
These are safe whenever, as long as they don't introduce new
evaluator primitives or change semantics:

\begin{itemize}
\item Pure stdlib libraries (\code{lib\_array}, list/string utilities)
\item REPL polish, CLI/help, docs, examples
\item Benchmark harness, profiling surfaces
\item Pretty-printing, parser ergonomics
\item Test cleanup, expected-file maintenance
\item \textbf{Scout libraries for the ASM seam} (anytime after
  Phase~2 begins): small pure-fallback libraries whose workloads
  empirically stress the episode seam without committing to broader
  library design---\code{lib\_set.metta} (\code{set-insert-new?},
  \code{set-contains?}), \code{lib\_queue.metta},
  \code{lib\_key.metta} (compact-key encodings, e.g.\ 9-tile board
  $\rightarrow$ 64-bit int).  These are seam probes, not an
  application/library unleashing.  Positive example: \code{lib\_set}
  plus a \code{board-key} helper driving the tilepuzzle dedup matrix.
  Negative example: \code{lib\_http}, sockets, or anything forcing
  new backend contracts before the seam is stable.
\end{itemize}

Rule of thumb: if it lives \emph{above} the evaluator contract and
doesn't require new language semantics, it's anytime work.

Semi-anytime (can start but should merge after relevant seam stabilizes):
Lean proofs of already-named invariants, backend adapters, Python polish,
\code{--lang petta} mode, advanced benchmark suites.

\section{What This Guide Doesn't Do}

It doesn't prescribe exact steps, timelines, or success metrics.
It doesn't define rollback protocols or kill criteria.  It doesn't
pretend to know more than the roadmap about where the clouds of
possibility lead.

The roadmap describes the territory.  This guide says: if you're
going to walk it, universal term sharing is the natural first step,
and here's why the rest follows.

Trust the work.  Let the workloads guide refinement.

\end{document}
